% Part: second-order-logic
% Chapter: syntax-and-semantics
% Section: semantic-notions

\documentclass[../../../include/open-logic-section]{subfiles}

\begin{document}

\olfileid[fr]{sol}{syn}{sem}
\olsection{Notions sémantiques}

\begin{explain}
Comme au premier ordre, la relation de satisfaction permet de définir
la \emph{validité}, la \emph{conséquence logique} et la
\emph{satisfaisabilité}. On emploie ici la relation de satisfaction
des !!{formula}s du second ordre. Un !!{sentence} du second ordre
est une !!{formula} sans occurrence libre d'!!a{variable}, qu'il
s'agisse d'une variable individuelle, de prédicat ou de fonction.
\end{explain}

\begin{defn}[Validité]
Un énoncé~$!A$ est \emph{valide}, ce que nous
notons $\Entails !A$, si et seulement si
$\Sat{M}{!A}$ pour toute
!!{structure}~$\Struct M$.
\end{defn}

\begin{defn}[Conséquence logique]
Un ensemble d'énoncés~$\Gamma$ \emph{entraîne}
un énoncé~$!A$, ce que nous notons
$\Gamma\Entails !A$, si et seulement si,
pour toute !!{structure}~$\Struct M$ telle que
$\Sat{M}{\Gamma}$, on a $\Sat{M}{!A}$.
\end{defn}

\begin{defn}[Satisfaisabilité]
Un ensemble d'énoncés~$\Gamma$ est
\emph{satisfaisable} s'il existe une
!!{structure}~$\Struct M$ telle que
$\Sat{M}{\Gamma}$. S'il n'est pas satisfaisable,
il est dit \emph{insatisfaisable}.
\end{defn}

\end{document}
